%@ pages 17-35
%@ status translated
% Original section 2, "Les n-catégories de Segal", pp. 17--35.
%
% Translation notes (the mathematics is left exactly as printed):
%  * "n-précat (de Segal)" is rendered "Segal n-precategory" / "n-precategory";
%    the categories nSePC, nPC, nSePCh and the construction SeCat keep their
%    printed names.
%  * "cmf" is rendered "cmc", per the README.
%  * "catégorie de modèles fermée engendrée par cofibrations" is rendered
%    "cofibrantly generated closed model category", the standard English term.
%  * "flèche" is "arrow", "k-groupique" is "k-groupic".
%  * In the proof of the lemma labelled secatunique the source's array diagram
%    has its second row identical to its first (F^1(A) -> F^1(F^2(A)) twice,
%    although the vertical arrow is labelled F^1(f) and F^1(F^2(f))). The
%    diagram is reproduced as printed; nothing is emended.

\clearpage
\section{Segal \texorpdfstring{$n$}{n}-categories}
\label{ncatsegpage}

Let us begin by fixing some notation: the categories of sets, of simplicial
sets, and of topological spaces are denoted respectively
$Ens$, $EnsSpl$ and $Top$. A simplicial set is a functor
$\Delta ^o\rightarrow Ens$ where $\Delta$ is the category whose objects are
the ordered sets $p:= \{ 0', \ldots , p'\}$ and whose morphisms are the
increasing maps. If $X$ is a simplicial set we shall denote by
$X_p$ its value on $p\in \Delta$. If $A: \Delta
^o\rightarrow \Cc$ is a functor, we shall denote by $A_p$ or sometimes $A_{p/}$
the value it takes on $p\in \Delta ^o$.
We shall denote by $|A|$ the topological realisation of the
simplicial set $A$. If $Y$ is a category (i.e.\ a $1$-category) we shall denote
its nerve by $\nu Y \in EnsSpl$. We shall use the abbreviation
``cmc'' for ``closed model category'' of \cite{Quillen}.

The following definitions are generalisations of definitions of
Tamsamani \cite{Tamsamani} and Dunn \cite{Dunn}, which in turn took up
ideas of Segal \cite{SegalTopology} \cite{Adams}.
%%%%%%%
The notion of Segal $1$-category appears for the first time
at the end of Dwyer-Kan-Smith \cite{DKS}.
%%%%%%%% above phrase added in v3 &&&&&&&&
See the
comparison below.

We define the notion of {\em Segal $n$-precategory}, and the category $nSePC$
of these objects, by recursion on $n$. A set may be regarded as a Segal
$n$-precategory: there will be a (fully faithful) functor
$Ens\rightarrow nSePC$.
An object in the image of this functor will be called ``a discrete set''.
For $n=0$, a {\em Segal $0$-precategory} is by definition a simplicial set;
and the functor $Ens\rightarrow 0SePC$ associates to a set $S$ the functor
$\underline{S}:\Delta ^o\rightarrow Ens$ constant with values $S$. For
$n\geq 1$, a {\em Segal $n$-precategory} is a functor
$$
\Delta ^o\rightarrow (n-1)SePC
$$
denoted $p\mapsto A_{p/}$, such that $A_{0/}$ is a discrete set which we
also denote
$A_0$. A {\em morphism} of Segal $n$-precategories is a natural
transformation of functors $\Delta ^o\rightarrow (n-1)SePC$, which defines the
category $nSePC$. The functor $Ens\rightarrow nSePC$ is the one which to a
set $S$ associates the composite functor
$$
\Delta ^o\stackrel{\underline{S}}{\rightarrow} Ens \rightarrow (n-1)SePC.
$$

One may unravel the recursion in this definition. If one wants to follow the
notation of Tamsamani \cite{Tamsamani}, a Segal $n$-precategory is therefore an
$(n+1)$-simplicial set, i.e.\ a functor
$$
(\Delta ^o)^{n+1}\rightarrow Ens
$$
denoted $(m_1,\ldots , m_{n+1})\mapsto A_{m_1,\ldots , m_{n+1}}$,
which satisfies the {\em constancy conditions} according to which,
for $m_i=0$ the
functor
is constant in the variables $m_{i+1},\ldots , m_{n+1}$.

These constancy conditions can be taken into account once and for all by
introducing the category $\Theta ^{n+1}$, a quotient of $\Delta ^{n+1}$
\cite{nCAT}. We recall that $\Theta ^{n+1}$ is the category whose objects are
the sequences $(m_1,\ldots , m_i)$ for $i\leq n+1$, with $m_i\in \Delta$ and
with the identifications
$$
(m_1,\ldots , m_j=0, \ldots , m_i) =: (m_1,\ldots , m_{j-1}).
$$
The object $(0,\ldots , 0)$ is also equivalent to the sequence of length
$0$ which is denoted $0$. The morphisms come from the morphisms of
$\Delta$ with the change induced by these identifications; cf.\ \cite{nCAT}
or \cite{limits}.

For $k\leq n$ and for $M=(m_1,\ldots , m_k)\in \Delta ^k$ we denote by $A_{M/}$
the Segal $(n-k)$-precategory
$$
(m'_1,\ldots , m'_{n+1-k})\mapsto A_{m_1,\ldots , m_k, m'_1,\ldots ,
m'_{n+1-k}}.
$$
For $k=1$ this agrees with the notation $A_{p/}$ introduced above.

A Segal $n$-precategory may be viewed as a functor
$$
(\Theta ^{n+1})^o\rightarrow Ens,
$$
with no further condition (the constancy conditions being taken into account by
the fact that $\Theta ^{n+1}$ is the adequate quotient of $\Delta ^{n+1}$).
As a result, $nSePC$ appears as the category of presheaves of sets on
$\Theta ^{n+1}$, and therefore admits all limits and colimits with small
indexing, as well as an internal $\underline{Hom}$.

In fact, the most useful point of view is a mixture of the two points of view
of the preceding paragraph: a Segal $n$-precategory may (indeed must) be
regarded as a functor
$$
A: (\Theta ^n)^o\rightarrow EnsSpl
$$
denoted $M\mapsto A_M$, such that, for $|M| < n$, the simplicial set $A_M$
is a discrete set. It is sometimes useful to confuse simplicial sets with
topological spaces, and one obtains essentially the same notion by considering
the functors $(\Theta ^n)^o\rightarrow Top$.

We shall say that $A_0$ is {\em the set of objects of $A$}. For two objects
$x,y\in A_0$
we shall denote by $A_{1/}(x,y)$ the preimage of $(x,y)\in A_0\times A_0$ under
the morphism
$$
A_{1/}\rightarrow A_0\times A_0
$$
induced by the two ``face'' maps of the simplicial structure.
We shall say that $A_{1/}(x,y)$ is {\em the Segal $(n-1)$-precategory of
morphisms between $x$
and $y$}. We shall sometimes use this same notation (which is quite compact)
for the set of morphisms between two objects of a category or for the
simplicial set of morphisms between two objects of a simplicial category.


We now define the {\em
Segal $n$-categories}, which are the Segal $n$-precategories satisfying
certain conditions of ``speciality'' generalising the original condition
of Segal \cite{Segal}. To formulate these conditions one needs a large
recursion as in \cite{Tamsamani}, cf.\ also the summary in
\cite{limits}. We shall define simultaneously by recursion on $n$: the
condition for a Segal $n$-precategory to be a Segal $n$-category; the
condition for a morphism between two Segal $n$-categories to be an
equivalence; and the ``truncated'' set
$\tau _{\leq 0}(A)$ associated to a
Segal $n$-category $A$. First for $n=0$, and by definition:
every Segal $0$-precategory, i.e.\
simplicial set, is a Segal $0$-category; a morphism is an
equivalence if it is a weak equivalence of simplicial sets, i.e.\
if it induces isomorphisms between the homotopy groups (and the $\pi
_0$) of the
topological realisations; and the truncation $\tau _{\leq 0}(A)$ of a
simplicial set $A$ is the set $\pi _0(|A|)$. For the rest, the recursion is
identical to Tamsamani's \cite{Tamsamani}: let $n\geq 1$ and suppose
that the definitions are available at level $n-1$. Let $A$ be a
Segal $n$-precategory, viewed as a functor
$$
A: \Delta ^o\rightarrow (n-1)SePC,\;\;\; p\mapsto A_{p/}.
$$
We shall say that $A$ is a {\em Segal $n$-category} if the
$A_{p/}$ are Segal $(n-1)$-categories, and if for every $p$ the morphism
of Segal $(n-1)$-precategories
$$
A_{p/}\rightarrow A_{1/}\times _{A_0} \ldots \times _{A_0}A_{1/}
$$
is an equivalence of Segal $(n-1)$-categories. This last
``Segal morphism'' being the one whose components are the maps
dual to the inclusions $\{ i,i+1\} \subset \{ 0,\ldots , p\}$ in $\Delta$,
see \cite{SegalInventiones}, \cite{Adams}, \cite{Dunn}, \cite{Tamsamani}.


If $A$ is a Segal $n$-category, then the functor
$$
p\mapsto \tau _{\leq 0} (A_{p/}) \;\;\;\; \Delta ^o \rightarrow Ens
$$
is the nerve of a category which we denote $\tau _{\leq 1}(A)$. We
define $\tau _{\leq 0}(A)$ as the set of isomorphism classes of the
category $\tau _{\leq 1}(A)$. If $A$ and $B$ are Segal $n$-categories
and $f:A\rightarrow B$ is a morphism (i.e.\ a morphism of $nSePC$) we shall
say that
$f$ is an {\em equivalence} if the induced morphism
$$
\tau _{\leq 0}(A)\rightarrow \tau _{\leq 0}(B)
$$
is surjective (we shall then say that ``$f$ is essentially surjective''), and
if for
every pair of objects $x,y\in A_0$, the morphism
$$
A_{1/}(x,y)\rightarrow B_{1/}(f(x), f(y))
$$
is an equivalence of Segal $(n-1)$-categories (we shall then say that
``$f$ is
fully faithful''). Compare with the notion
of equivalence between simplicial categories of
\cite{DwyerKanDiags}.


It is only for a Segal $n$-category $A$ that it is reasonable to
consider
the $(n-1)$-category of morphisms between two objects $A_{1/}(x,y)$.
A {\em $1$-arrow of $A$} is an object of $A_{1/}(x,y)$.
By iteration one obtains the notion of $i$-arrow of $A$ \cite{Tamsamani}.
Let $1^i\in \Theta ^n$ be the object $(1, \ldots , 1)$ ($i$ times), image of
the object
$(1, \ldots , 1, 0,\ldots , 0)\in \Delta ^n$. An {\em $i$-arrow of $A$}
($i<n$) is an element of the set $A_{1^i}$. For $i=n$, $A_{1^n}$ is a
simplicial set and an {\em $n$-arrow of $A$} is a vertex (i.e.\
an element of the degree $0$ part) of the simplicial set
$A_{1^n}$. The face morphisms in the $i$-th simplicial variable
provide the {\em source} and {\em target} maps which to an $i$-arrow
associate $(i-1)$-arrows (a $0$-arrow is an object, i.e.\ an element
of $A_0$).

Let $A$ be a Segal $n$-category.
Following \cite{Tamsamani}, we shall say that an $i$-arrow ($i\geq 1$) is {\em
invertible up to equivalence} or {\em up to homotopy} if its image
as a $1$-arrow in the $1$-category $\tau _{\leq 1}(A_{1^{i-1}/})$ is
invertible. See \cite{Tamsamani} for more details. We shall say (still
following \cite{Tamsamani}) that $A$ is a {\em Segal $n$-groupoid}
if for
every $1\leq i\leq n$, all the $i$-arrows are invertible up to
equivalence. More generally, for $0\leq k \leq n$ we shall say that $A$ is
{\em $k$-groupic} if the $i$-arrows are invertible up to
equivalence for every $i>k$. In particular, $A$ is always $n$-groupic, and
$A$ is $0$-groupic if and only if it is a Segal $n$-groupoid.

\subnumero{Comparison with Tamsamani and Dunn}
We compare with \cite{Tamsamani} and \cite{Dunn}. Tamsamani
(\cite{Tamsamani}, 1995)
considers the functors $(\Delta ^o)^n\rightarrow Ens$ and imposes on them the
constancy conditions as well as the Segal conditions cf.\ below,
obtaining the notion of {\em (weak) $n$-category}. In \cite{nCAT},
one introduces the notion of {\em $n$-precategory} by relaxing all but the
constancy conditions.
An $n$-precategory (resp.\ $n$-category of Tamsamani),
which is a functor $\Theta ^n\rightarrow Ens$, may be viewed as a
Segal $n$-precategory by composing with the inclusion $Ens\rightarrow EnsSpl$.


More generally, the Segal $n$-precategories
(resp.\ Segal $n$-categories)
such that the values of the functor $\Theta ^n \rightarrow EnsSpl$
are $0$-truncated simplicial sets (i.e.\ disjoint unions of
contractible components) correspond to $n$-precategories
(resp.\ $n$-categories) by composition with $\pi _0: EnsSpl
\rightarrow Ens$. With this translation, most of our statements on
Segal $n$-categories remain true {\em mutatis mutandis} for
$n$-categories (the only exception concerns the descent
property---not preserved by truncation---for Dwyer-Kan localisations,
which are not truncated in general). We often formulate
our
statements only in the framework of Segal $n$-categories, leaving it to the
reader to write the corresponding statements for the
$n$-categories.
On the other hand, we shall return later to the compatibilities between the
notions of
$n$-category and Segal $n$-category (and also to the compatibilities
between different values of $n$).

Dunn, in \cite{Dunn} (submitted in 1993, published in 1996 but summarising his
1984 thesis, and which refers to Cobb \cite{Cobb})
considers---among other ``delooping machines''---the {\em $n$-th
iterate of the Segal machine} whose underlying category is that of the
functors
$$
A: (\Delta ^o) ^n\rightarrow Top
$$
such that $A_{m_1,\ldots , m_n}=\ast$ if $m_i=0$ (cf.\ \cite{Dunn}, beginning of
\S 3). One may call these objects ``Dunn $n$-precategories''.
Replacing $Top$ by $EnsSpl$ and passing to the quotient, one obtains
the Segal $n$-precategories $A$ satisfying $A_M=\ast$ for
$|M|<n$.

Dunn also imposes on his objects the speciality condition in each simplicial
variable
(\cite{Dunn} Definition 3.1) and one could call the objects which satisfy
this condition the ``Dunn $n$-categories''. One can check that
this
condition coincides with the condition of being a Segal $n$-category, in
the case he considers where $A_M= \ast$ for $|M|< n$.
In sum, Dunn had looked at the Segal $n$-categories having a
single $i$-arrow for every $i< n$.

The definitions of Dunn and of Tamsamani
are based on the idea of iterating the construction of Segal
\cite{Segal}. They are complementary as the
Postnikov tower and the Whitehead tower are: with Dunn (as with Whitehead)
it is a question of homotopy in degrees $\geq n$ whereas with Tamsamani (as
with Postnikov)
it is a question of homotopy in degrees $\leq n$. The notion of {\em Segal
$n$-precategory} makes it possible to integrate these two points of view.
Moreover, in the ``$n$-categorical'' variables one can have arrows which are
not (homotopically) invertible. It is therefore a definition which allows one
to have homotopy in every degree, and non-invertible arrows in degrees $\leq
n$. It is an avatar of the notion (not yet fully elucidated) of {\em
$\infty$-precategory} which
would allow one to have non-invertible arrows in every degree. We have
decided to do the present writing in the framework of Segal $n$-precategories,
for the sake of simplicity and also because, in the examples which we have in
view,
the non-invertibility of arrows comes into play only for a finite number
of degrees.


\subnumero{The \texorpdfstring{$SeCat$}{SeCat} operation}

One of the basic ideas which made it possible to obtain the structure of
closed model category for $n$-precategories in \cite{nCAT} was
the operation $Cat$, which one may sum up very simply by saying that it forces
the condition of being an $n$-category. This procedure applies in the same
way in
our framework of Segal $n$-precategories: we shall have an operation which we
denote
$A\mapsto SeCat(A)$ with a natural transformation $i_A: A\rightarrow
SeCat(A)$, which will transform every Segal $n$-precategory into a
Segal $n$-category.

In fact, instead of fixing the operation $SeCat$ precisely, we
give the following definition.

\begin{definition}
\label{secat}
One calls {\em operation of type $SeCat$} any pair $(F, i)$ where
$F:nSePC\rightarrow nSePC$ is a functor and $i_A: A\rightarrow F(A)$ a
natural transformation, possessing the following properties:
\newline
(a)---for every Segal $n$-precategory $A$, $F(A)$ is a Segal
$n$-category;
\newline
(b)---$i_A$ is an isomorphism on the sets of objects, and if $A$ is a
Segal $n$-category then $i_A$ is an equivalence of Segal
$n$-categories; and
\newline
(c)---for every Segal $n$-precategory $A$, the morphism $F(i_A)$ is an
equivalence of Segal $n$-categories.
\end{definition}

\begin{lemma}
\label{secatunique}
There exists an operation of type $SeCat$. If $(F^1, i^1)$ and $(F^2, i^2)$
are two operations of type $SeCat$ then they are equivalent in the
sense that
$$
F^1(i^2_A): F^1(A) \rightarrow F^1(F^2(A))
$$
is an equivalence of Segal $n$-categories for every Segal $n$-precategory
$A$. In particular, if $f:A\rightarrow B$ is a morphism of
Segal $n$-precategories then $F^1(f)$ is an equivalence of Segal
$n$-categories if and only if $F^2(f)$ is.
\end{lemma}
{\em Proof:}
For existence, one may use a definition analogous to that of
\cite{nCAT}, which treats the case of (non-Segal) $n$-precategories (for
another approach, see the discussion
further on).
For the
simplicial sets of the last degree,
one must fill all the ``horns'' (anodyne Kan extensions). For uniqueness,
one uses the same proof as that of
\cite{nCAT} Proposition 4.2. In fact this proof must be coupled
with the proof of theorem \ref{cmf} in a recursion on $n$
as is done in \cite{nCAT}. We note that the case $n=0$ is trivially
true.

If $f: A\rightarrow B$ and if $F^2(f)$ is an equivalence then in the
diagram
$$
\begin{array}{ccc}
F^1(A) & \stackrel{F^1(i^2_A)}{\rightarrow} &  F^1(F^2(A))\\
{\scriptstyle F^1(f)}\downarrow && \downarrow {\scriptstyle F^1(F^2(f))}\\
F^1(A) & \stackrel{F^1(i^2_A)}{\rightarrow} &  F^1(F^2(A))
\end{array}
$$
the arrow $F^1(F^2(f))$ is an equivalence because $F^1$ transforms
equivalences of Segal $n$-categories into equivalences by
\ref{secat} (b). The horizontal arrows are equivalences by the
uniqueness result above. Hence $F^1(f)$ is an equivalence.
\eop


We shall call $SeCat$ any operation which satisfies the criteria of
definition \ref{secat}. We now give a convenient variant, already
present (for $n$-precategories) in \cite{nCAT} and explored
in \cite{effective} for
Segal $1$-precategories. For $n=0$ one may take the identity
functor with the identity natural transformation. Let $n\geq 1$
and suppose the operation $SeCat$ constructed for Segal $(n-1)$-precategories.
Let $A$ be a Segal $n$-precategory. We define a Segal $n$-precategory
$Fix(A)$
with the same objects as $A$ by setting (for $p\geq 1$) $$
Fix(A)_{p/}(x_0,\ldots , x_p):= SeCat(A_{p/}(x_0,\ldots , x_p)).
$$
This has the effect of transforming the $A_{p/}$ into Segal $n$-categories.
We have a natural morphism $A\rightarrow Fix(A)$.
On the other hand, for every $m\geq 2$ we define a Segal $n$-precategory
$Gen[m](A)$ equipped with a morphism $A\rightarrow Gen[m](A)$ inducing an
isomorphism on the sets of objects, by
$$
Gen[m](A)_{q/} := A_{q/} \cup ^{\coprod_{q\rightarrow m} A_{m/}}
\coprod _{q\rightarrow m}\Gg [m]
$$
where
$$
A_{m/}\stackrel{a}{\rightarrow} \Gg [m]
\stackrel{b}{\rightarrow} A_{1/}\times _{A_0}\ldots
\times _{A_0} A_{1/}
$$
is a factorisation of the Segal morphism, where $a$ is a
cofibration and $b$ a weak equivalence, (which exists by the theorem
\ref{SeCmf} which we suppose known for $n-1$). The coproducts in the definition
of $Gen[m]$ are taken with all the morphisms $q\rightarrow m$ in
$\Delta$ which
do not
factor through a principal edge $1\rightarrow m$. See
\cite{nCAT} or \cite{effective} for how to define
the functoriality morphisms $Gen[m](A)_{q/}\rightarrow
Gen[m](A)_{p/}$ for $p\rightarrow q$.
Now we define $Gen$ as the composite for every $m\geq 2$ of the
$Gen[m]$ and $SeCat$ is obtained by iterating a countably infinite
number $\omega$ of times the composite $Fix \circ Gen$.
This operation, which appears very ineffective, in reality has quite
good effectivity properties, see \cite{effective}.

One may view an $n$-precategory $A$ as a {\em system of generators
and relations} for defining an $n$-category, and $SeCat(A)$ as the
$n$-category thus generated.

\subnumero{The closed model category structure}
We now define the closed model category structure on
$nSePC$.

A morphism $A\rightarrow B$ of Segal $n$-precategories will be a {\em
cofibration} if $A_M\rightarrow B_M$ is injective for every $M\in \Theta ^{n+1}$
(this condition is different from the condition for the
$(n+1)$-precategories of \cite{nCAT} because here we also require
injectivity for $|M|=n+1$). A cofibration is
therefore simply an injection of presheaves
of sets.

We shall say that a morphism $A\rightarrow B$ of Segal $n$-precategories is
a {\em
weak equivalence} if the morphism
$SeCat(A)\rightarrow SeCat(B)$ is an equivalence
of Segal $n$-categories. This property is independent of the choice of
the construction $SeCat$ thanks to lemma \ref{secatunique}.

And we shall say that a morphism $A\rightarrow B$ of Segal $n$-precategories
is a {\em
fibration} if it possesses the lifting property with respect to
trivial cofibrations (those which are weak equivalences). Let us
recall (\cite{Quillen}) that, by definition,
the lifting property ensures that
for every trivial cofibration $E\rightarrow E'$, and every pair of morphisms
$E\rightarrow A$ and $E'\rightarrow B$ which form a commutative square, there
exists a morphism $E'\rightarrow A$ making commutative the two triangles
in the square.

\begin{theorem}
\label{SeCmf}
The three classes of morphisms
above
make the category of Segal $n$-precategories
a closed model category.
\end{theorem}

The proof enters the more general framework of the comparison
between the notions of Segal $n$-category and/or $n$-category, for
different values of $n$. We broach this discussion before
the proof of \ref{SeCmf}.

\subnumero{Changing \texorpdfstring{$n$}{n}}

The closed model category of Segal $n$-precategories of
\ref{SeCmf} is a kind of substitute for the notion of $\infty$-category,
in which the $i$-arrows are invertible up to equivalence, for
$i>n$. If one admits the existence of a suitable theory of $\infty$-categories,
the equivalence
will pass through the construction $A\mapsto \Pi _{\infty} \circ A$ which
transforms a
Segal $n$-precategory into an $\infty$-precategory, with a $\Pi _{\infty}$
which makes correspond an
$\infty$-groupoid to every topological space or Kan simplicial set.
We do not wish to enter into the details of the notion of
$\infty$-category
here, because the version ``Segal $n$-categories'' suits our
applications.
This choice is all the more reasonable in that one can construct
functors $A\mapsto \Pi
_{m,Se}\circ A$ to pass from $n$ to $n+m$.




Let us recall that Tamsamani \cite{Tamsamani} has defined a functor
``Poincaré
$n$-groupoid''
$$
\Pi _n : Top \rightarrow nCat
$$
where $nCat$ is the full subcategory of the category $nPC$ of
$n$-precategories, formed of the $n$-categories. This functor generalises
the classical ``Poincaré groupoid'' $\Pi _1(X)$. Tamsamani has shown that
this functor establishes an equivalence between the homotopy theories
of $n$-truncated topological spaces (i.e.\ where the $\pi _i$ vanish for
$i>n$) and of $n$-groupoids (i.e.\ $n$-categories where the arrows are
invertible up to equivalence). More precisely, he constructed a
``topological realisation'' functor
$$
\Re : nCat \rightarrow Top
$$
and equivalences $\Pi _n \circ \Re\cong 1$ and $\Re \circ \Pi _n \cong 1$
for the restrictions of these functors to the categories $Top _n$ of
$n$-truncated topological spaces and $nGpd$ of $n$-groupoids. We shall
therefore adopt the
point of view that $n$-groupoids may be freely replaced by
$n$-truncated topological spaces (or simplicial sets).

We make here some remarks which use notions that will be
treated later in the present work, in order to offer another light
on this equivalence result of Tamsamani (however,
the reader may do without it until the symbol ``$\oslash$'').
Applying the
Dwyer-Kan theory \cite{DwyerKan3} which will be recalled in \S 8
below, one can invert the equivalences and obtain
``localised'' simplicial categories
$L(Top_n)$ and  $L(nGpd)$.
Now the functors $\Pi _n$ and $\Re$ and
Tamsamani's equivalences, in conjunction with proposition 3.5 and
corollary 3.6 of \cite{DwyerKan2} (cf.\ also our lemma \ref{nonadjoint}),
give an equivalence
$$
L(Top_n)\cong L(nGpd )
$$
of simplicial categories. One can push this a little further by using
our notations of \S 11: we denote by $nCAT$ the $(n+1)$-category of
$n$-categories (see \cite{nCAT}).
\footnote{
The reader will take care to distinguish between the typography in capitals
$nCAT$
and the typography which involves lower-case letters $nCat$.
Indeed, $nCAT$ denotes the $(n+1)$-category of $n$-categories,
defined with the help of the internal structure cf.\ \S 11 and (\cite{nCAT}
\S 7); whereas
$nCat$ denotes
the strict $1$-category with the same objects and the morphisms which strictly
respect
the structure (in particular, it is the latter that Tamsamani had
considered). One uses this same convention for stacks further on.}
%
%
%
Let us denote by $nGPD\subset nCAT$ the full sub-$(n+1)$-category of
$n$-groupoids. Then $nGPD$ is $1$-groupic (i.e.\ its
$n$-categories of
$Hom$ are in fact $n$-groupoids) hence equal to its $1$-groupic interior:
$nGPD^{int, 1}= nGPD$. By theorem \ref{intereqloc} one has
$$
L(nPC)\cong L(nPC_f)
\cong nCAT ^{int, 1}.
$$
On the other hand, since we have functorial fibrant replacements, every
category comprised between $nPC$ and $nPC_f$ gives the same Dwyer-Kan
localisation; this applies notably to $nPC_f \subset nCat \subset nPC$.
One thus obtains the equivalence of $(n+1)$-categories
$$
L(nCat)\cong nCAT^{int, 1}
$$
(more technically, this is here an equivalence which possibly passes through
morphisms in both directions, since neither of the two $(n+1)$-precategories in
question is
fibrant).
The subcategory $nGpd\subset nCat$ is defined by a condition
invariant under equivalence, hence (by \ref{stabilite} below) $L(nGpd)$
is the full subcategory of $L(nCat)$ formed of the objects which are
$n$-groupoids. By definition the same is true for $nGPD$ and hence
also for $nGPD^{int, 1}$. One thus obtains the equivalence
$$
L(nGpd) \cong nGPD.
$$
If one now applies the (immediate) variant
of Tamsamani's comparison theorem
which gives the equivalence between Dwyer-Kan localisations,
one obtains an equivalence of $(n+1)$-categories
$$
L(Top_n)\cong nGPD.
$$
On the other hand, one can observe that the
simplicial sets
$L(Top _n)_{1/}(X,Y)$ have the right homotopy type, namely that
of the mapping space
$\underline{Hom}(X,Y)$, as soon as, for example, $X$ is a CW-complex
(see \cite{DwyerKan3} etc.).
To put
things differently, $L(Top_n)$ is equivalent to the simplicial category
of
($n$-truncated) fibrant and cofibrant objects for any simplicial closed model
category structure for topological spaces. Hence
$L(Top_n)$ is ``the'' $(n+1)$-category of $n$-truncated spaces. This
equivalence therefore gives precisions on what one means by
``$n$-groupoids are identified with $n$-truncated topological spaces''.

\noindent $\oslash$

\medskip

In view of the preceding remarks (the short version would have sufficed), we
are going to
{\em declare} that an $\infty$-groupoid is nothing other than a topological
space or simplicial set. This point of view, which already guides us
in the
definition of Segal $n$-categories, also guides us in
comparing these notions for different values of $n$.


As explained at the beginning of this chapter, one may view every Segal
$n$-precategory as a simplicial presheaf over $\Theta ^n$, i.e.\ a
functor to the category of simplicial sets,
$$
A:\Theta ^n \rightarrow EnsSpl,\;\;\;\; M\mapsto A_M
$$
such that for $|M| < n$ the simplicial set
$A_M$ is a discrete set.

From this point of view, one may
compose with the functor
$| \; |$ of topological realisation of simplicial sets, then with
Tamsamani's functor $\Pi _m$ \cite{Tamsamani}. Equivalently
one may first apply Kan's functor
$Ex^{\infty}$ to arrive at Kan simplicial sets, then the variant
of the functor $\Pi _m$ of \cite{Tamsamani} adapted to simplicial
sets.  By these two (essentially equivalent) procedures, one obtains an
$(n+m)$-precategory which we denote  $\Pi _m \circ A$.

In the other direction, if one has an $(n+m)$-precategory $B$, one may apply to
it
the operation $\Re _{\geq n}$ which consists in taking the topological
realisation (or rather as a simplicial set) of \cite{Tamsamani} of the
$m$-precategories $B_{M/}$ for $|M|=n$ (one keeps the discrete sets $B_M$ for
$|M|<n$). The result is a Segal $n$-precategory
$\Re _{\geq n}(B)$. If $A$ is a Segal $n$-category having only
$m$-truncated simplicial sets in the image of the functor
$$
(M\mapsto A_{M/}): (\Theta ^n )^o\rightarrow EnsSpl
$$
considered above, then $\Re _{\geq n}\Pi _m\circ A$ is
equivalent to
$A$. If $A$
does not satisfy this truncation condition then this construction has
the effect
of performing the Postnikov truncation on the simplicial sets.
On the other hand,
if $A$ is an $n$-groupic $(n+m)$-category (i.e.\ where the $i$-arrows are
invertible up to homotopy for $i>n$), then $\Pi _m\circ \Re _{\geq
n}(A)$
is equivalent to $A$. In the opposite case one thus obtains the
{\em $n$-groupic completion of $A$} (i.e.\ the universal morphism
from $A$ to an $n$-groupic $(n+m)$-category),
analogous to the ``group-completion'' of a
simplicial monoid in algebraic topology (which is the case $n=0$).

The construction $A\mapsto \Pi _{m}\circ A$ is compatible up to homotopy
with coproducts. If $A\rightarrow B$ is a cofibration (i.e.\ injection,
cf.\ the definition below) and $A\rightarrow C$ a morphism of Segal
$n$-precategories, the coproduct $B\cup ^AC$ is a coproduct of simplicial sets
object by object over $\Theta ^n$. One may therefore apply the
Seifert-Van Kampen theorem for $\Pi _m$ of \cite{nCAT}, object by object over
$\Theta ^n$, to obtain the equivalence
$$
(\Pi _m \circ B)\cup ^{(\Pi _m \circ A)}(\Pi _m \circ C)
\stackrel{\cong}{\rightarrow}
\Pi _m \circ (B\cup ^AC).
$$
In the same order of ideas one remarks that if $A$ is a Segal $n$-precategory
then there is a natural equivalence of Segal $m$-categories
$$
\Pi _m \circ SeCat(A) \stackrel{\cong}{\rightarrow} Cat(\Pi _m \circ A).
$$



One may modify the
paragraph on the operation $\Pi _m\circ A$ to obtain a
passage from Segal $n$-categories to Segal $(n+m)$-categories.
For this, one must
first define a $\Pi _{m,Se}(X)$ for every topological space $X$,
which will be a Segal $m$-category.  For $m=0$ it is quite simply the
simplicial set ``singular complex'' $Sing(X)$. For general $m$ we
first apply the functor $\underline{\Omega}$ defined by
Tamsamani
\cite{Tamsamani}, $m$
times iterated as in
\cite{Tamsamani}, which gives a functor
$$
\underline{\Omega}^m(X): \Theta ^m\rightarrow Top.
$$
Then we compose with the functor $Sing$ to obtain
$$
\Pi _{m,Se}(X):= Sing \, \underline{\Omega}^m(X).
$$
Now if $A$
is a Segal $n$-precategory one considers it as
a functor
$$
\Theta ^n\rightarrow EnsSpl
$$
which one composes first with the realisation $| \; |$, then with the functor
$\Pi _{m,Se}$. We shall denote  $\Pi _{m,Se}\circ A$ the result of
this construction. It is a
Segal $(n+m)$-precategory. Likewise if $B$ is a Segal $(n+m)$-precategory
one may
define its {\em realisation in degrees $\geq n$} $\Re _{\geq n}(B)$.
One has the equivalence
$$
A \cong \Re _{\geq n}(\Pi_{m,Se}\circ A)
$$
(this time without truncation condition on $A$).
If moreover $B$ is $n$-groupic, then
one has $B\cong \Pi_{m,Se}\circ \Re _{\geq n}(B)$. And again, in the general
case, $\Pi_{m,Se}\circ \Re _{\geq n}(B)$
will be the $n$-groupic completion of $B$.

Finally one has the formula
$$
(\Pi _{m,Se} \circ B)\cup ^{(\Pi _{m,Se} \circ A)}(\Pi _{m,Se} \circ C)
\cong
\Pi _{m,Se} \circ (B\cup ^AC).
$$

\subnumero{Induction}

A (non-Segal) $n$-precategory $A$ may be regarded as
a Segal $m$-precategory for every $m\geq n$, which one may denote
$Ind _n^{m,Se}(A)$ with
$$
Ind _n^{m,Se}(A)_{k_1, \ldots , k_{m+1}}:= A_{k_1, \ldots , k_n}
$$
for $(k_1,\ldots , k_{m+1})\in \Delta ^{m+1}$.

{\em Warning:} If $A$ is an $n$-precategory then $SeCat(Ind _n^{m,Se}(A))$
is not in general equivalent to $Ind _n^{m,Se}(Cat(A))$. One has this
equivalence only if $A$ is already an $n$-category.


On the other hand it must be remarked that if  $A$ is a (non-Segal)
$n$-category
then $Ind _n^{m,Se}(A)$ is from the outset a Segal $m$-category.

For all these
reasons, one will use this ``induction'' operation
most often with an argument $A$ which is already an $n$-category.

The case which one will encounter most is that where $Y$ is a
$1$-category, which one
regards as a $1$-precategory by taking its nerve $\nu Y$. Then
for every $n\geq 1$ one has at one's disposal the Segal $n$-category
$Ind _n^{m,Se}(\nu Y)$.

In practice one will confuse $Y$ with the
induced Segal $n$-category.



\subnumero{Truncation}

In the framework of the above comparisons,
we now define a {\em truncation} $A\mapsto \tau _{\leq n}(A)$
for Segal
$m$-categories. The result is
a (non-Segal) $n$-category.
If $n\geq m$ we set
$$
\tau _{\leq n} (A) := \Pi _{n-m} \circ A
$$
with the notations introduced above. In particular one has
$$
\tau _{\leq m}(A)= \pi _0\circ A.
$$
For $n<m$ we set
$$
\tau _{\leq n}(A):= \tau _{\leq n}(\tau _{\leq m}(A)),
$$
where the truncation $\tau _{\leq n}$ on the right is the truncation
which transforms an $m$-category into an $n$-category (see \cite{Tamsamani}).
In terms of the preceding definitions, one has
$$
\tau _{\leq n}(A)_M = A_M \;\; \mbox{for}\;\; |M|<n, \;\;\;\;
\tau _{\leq n}(A)_M = \tau _{\leq 0}(A_{M/}) \;\; \mbox{for}\;\; |M|=n.
$$

This truncation is essentially the ``Postnikov truncation''. If one wants
to consider that our Segal $m$-categories $A$ are in fact
$\infty$-categories, one may prefer the notation $\tau _{\leq n}(A)$
to the notation $\Pi _{n-m}\circ A$.


\begin{lemma}
A morphism $A\rightarrow B$ of Segal $m$-categories is an equivalence
if and only if $\tau _{\leq n}A\rightarrow \tau _{\leq n}B$ is an
equivalence of $n$-categories for every $n$.
\end{lemma}
{\em Proof:}
This is true for $m=0$, i.e.\ for simplicial sets, and one obtains the general
case by an obvious recursion.
\eop

Let $A$ be a Segal $m$-category. We shall say that $A$ is
{\em contractible} if the canonical morphism $A\rightarrow \ast$ is an
equivalence. We shall say that $A$ is {\em $0$-truncated} if $A$ is
a disjoint union of contractible Segal $m$-categories (here the
union may possibly be empty). Let $n\leq m$; we shall say that
$A$ is {\em $n$-truncated} if for every
$M$ with $|M| = n$, the Segal $(m-n)$-category $A_{M/}$ is $0$-truncated.
Finally, let $n\geq m$; we shall say that $A$ is {\em $n$-truncated} if for
every
$M$ with $|M|=m$, the simplicial set $A_{M/}$ is $(n-m)$-truncated (i.e.\ its
homotopy groups vanish in every degree $i>n-m$).

For $n\geq m$ one has a natural morphism
$$
\Pi _{n-m, Se}\circ A\rightarrow
\tau _{\leq
n}(A).
$$
On the other hand, for $n<m$ there does not in general exist a morphism
$A\rightarrow Ind ^{m,Se}_n(\tau _{\leq n}(A))$ (nor in the other direction).
This results from
the possible existence of non-invertible $i$-arrows (the notion of interior
introduced below
will make it possible to get around this difficulty).  However, if $A$ is
$n$-truncated then one has a natural equivalence
$$
A\stackrel{\cong}{\rightarrow} Ind ^{m,Se}_n(\tau _{\leq n}(A)).
$$
In fact, one can check that a Segal $m$-category
$A$ is equivalent to one of the form $Ind ^{m,Se}_n(B)$,
where $B$ is an $n$-category, if and only if $A$ is $n$-truncated.




\subnumero{The interior}

The realisation provides a means of passing from a Segal $n$-precategory to a
simplicial set; this corresponds essentially to inverting all the
arrows in order to obtain an $\infty$-groupoid which one may view as a
simplicial set. This operation is rather brutal
and is hardly compatible
with questions of descent. We introduce here a dual construction (in a fairly
weak sense), {\em the interior} of a Segal $n$-precategory, which is much
better adapted to questions of descent and which will allow us
to unravel recursively the notion of stack in terms of a notion already well
known for simplicial presheaves.

If $A$ is a Segal $n$-category, we define {\em the interior of $A$}
denoted $A^{int, 0'}$ first as the sub-Segal-$n$-category of $A$
containing only the $i$-arrows invertible up to homotopy.
More precisely we shall say that an $i$-arrow $a\in A_{1^i}$ is {\em
invertible
up to homotopy} if the image of $a$ in $\tau _{\leq i}A$ is an
invertible $i$-arrow (let us recall that the composition of $i$-arrows in the
$i$-category $\tau _{\leq i}A$ is well defined, associative, etc.).
The interior $A^{int, 0'}$ of $A$ is the sub-Segal-$n$-precategory
defined by the
condition that $a\in A_M$ is in $A^{int,0'}_M$ if and only if
$f^{\ast}(a)$ is an $i$-arrow invertible up to homotopy, for
every morphism $f: 1^i\rightarrow M$ of $\Theta ^{n+1}$, for $i\leq n $.
One sees that $A^{int, 0'}$ is a Segal $n$-category, and
more precisely a
Segal $n$-groupoid.  We now introduce
$$
A^{int, 0}:= \Re _{\geq 0} A^{int, 0'},
$$
which is a simplicial set.
This last transformation does not change much because $A^{int, 0'}$
is already
a Segal $n$-groupoid, and in particular $\Pi _{n,Se}\circ A^{int,
0}$ is
equivalent to $A^{int, 0'}$.

In terms of $\infty$-categories,
$A^{int, 0}$ is the universal $\infty$-groupoid equipped with a morphism to
$A$.


More generally one defines {\em the $k$-groupic interior}
$A^{int, k}$ of a Segal $n$-category $A$, for every $k<n$. First one
defines $A^{int, k'}\subset A$ as the sub-Segal-$n$-precategory with
$a\in A_M$ contained in $A^{int, k'}_M$ if and only if $f^{\ast}a$
is an $i$-arrow invertible up to homotopy for every morphism $f:
1^i\rightarrow M$ with $k<i\leq n$. In other words, it is only
for $i>k$ that one keeps the
$i$-arrows invertible up to homotopy. Here again, one checks that
$A^{int, k'}$ is a $k$-groupic Segal $n$-category (i.e.\ whose
$i$-arrows are invertible up to homotopy for $i>k$). As
previously, one introduces
$$
A^{int, k}:= \Re _{\geq k} A^{int, k'},
$$
which is a Segal $k$-category. One has
$$
\Pi _{n-k,Se}\circ A^{int, k} \cong A^{int, k'}.
$$
In terms of $\infty$-categories,  $A^{int, k}$ is the
$k$-groupic $\infty$-category universal with a morphism to $A$.

If one wants to apply these constructions to a Segal $n$-precategory one must
first apply the operation $SeCat$ to obtain a Segal $n$-category
(where appropriate, this will be understood in the notation:
$A^{int, k}:= SeCat(A)^{int, k}$).

For $k<l<n$ one has $(A^{int, l})^{int, k}=A^{int , k}$ up to an
equivalence---which we shall omit to mention. On the other hand one has
$$
\tau _{\leq m}(A^{int, k})=\tau _{\leq m} A
$$
for $m\leq k$, and
$$
\tau _{\leq m} (A^{int, k})=(\tau _{\leq m})^{int, k}
$$
for $k<m$.
If $A\rightarrow B$ is an equivalence of Segal $n$-categories then so
is $A^{int, k}\rightarrow B^{int , k}$. Moreover, $A$ is
$k$-groupic if and only if $A^{int, k}\rightarrow A$ is an
equivalence.

Finally, one can use these constructions to relate a Segal $n$-category
$A$ to its truncation $\tau _{\leq k}(A)$. Let us recall that for $k\geq n$
one already has a morphism $\Pi _{k-n,Se}\circ A\rightarrow \tau _{\leq
k}(A)$. If moreover $k<n$, one has a pair of morphisms
$$
\Pi _{n-k,Se} \circ A^{int, k} \rightarrow A
$$
and
$$
A^{int, k} \rightarrow \tau _{\leq k}(A).
$$




\subnumero{How to obtain a closed model category structure}

Before broaching the proof of theorem \ref{SeCmf},
we extract from \cite{Jardine}, \cite{JardineGoerssBook},
\cite{Hirschhorn}, and more especially from Dwyer-Hirschhorn-Kan \cite{DHK} a
statement which expresses exactly what must be proved in order to obtain a
closed model category structure as in \ref{SeCmf}. A statement of the same
nature is found in Hovey's book (\cite{HoveyBook}
Theorem 2.1.11).


Let $M$ be a category and $C\subset M$ a subcategory (i.e.\ class of
morphisms). Let moreover $I\subset C$ be a subset of morphisms. One says
that {\em $I$ generates $C$} if, for every morphism $f:B\rightarrow B'$ of
$M$, the lifting property \cite{Quillen} with respect to the morphisms of
$I$ ensures the lifting property with respect to all the morphisms of
$C$  ---i.e.\ for every diagram
$$
\begin{array}{ccc}
E&\rightarrow & B \\
\downarrow && \downarrow \\
E'&\rightarrow & B'
\end{array}
$$
with $E\rightarrow E'$ in $C$,
the diagonal morphism exists.

We shall mention, without giving a definition, the property for $I$ of {\em
permitting the small object argument} (``small object argument''), cf.\
\cite{Hirschhorn}  \cite{JardineGoerssBook} and particularly \cite{DHK},
7.3. We note that if $M$ admits a suitable ``underlying set'' functor
$M\rightarrow Ens$, then every subset $I$ of morphisms
permits the small object argument (\cite{DHK} 8.5, \cite{Hirschhorn}). This
will be the case in all the examples we are going to consider. In fact, in
this
case one
will be able to say that every subcategory $C\subset M$ which is defined by
set-theoretically reasonable conditions will have a generating subset
$I\subset C$ permitting the small object argument---see for example
Jardine's argument \cite{Jardine}. In general we shall not
give a proof for this type of condition.

Let us now explain how we shall speak of ``lifting
property''.
There are two cases: either one is given the morphism $B\rightarrow B'$ and
one wants to have the lifting property in the above diagram for every
morphism $E\rightarrow E'$ in a certain class; or one is given
$E\rightarrow E'$ and one wants to have the property for every $B\rightarrow B'$
in a certain class. Instead of making the distinction by speaking of
``left lifting'' and ``right lifting'' (as Quillen does
\cite{Quillen}---but we cannot manage to remember where the left is in
this business)
we shall leave it to the reader to deduce from the context what is meant,
knowing that the ``cofibrations'' are always on the left and the
``fibrations'' on the right in this kind of diagram.


We recall the notion of {\em cofibrantly generated closed model category}
of Hirschhorn \cite{Hirschhorn} \cite{DHK}: it is a cmc $M$,
admitting all limits and colimits, and such that the
subcategories $cof\subset
M$ and $cof \cap W \subset M$ of cofibrations and of trivial cofibrations
admit generating sets permitting the small object argument.

We give here a version of the {\em recognition lemma}
(\cite{DHK} 8.1, which Hirschhorn attributes to Kan cf.\ \cite{Hirschhorn}). Our
version is not the most general, but on the other hand incorporates some
key remarks which are found only further on in \cite{DHK}, at the
moment of the proof of  the cmc structure for simplicial categories (\cite{DHK}
48.7). The lemma and its proof are contained implicitly
in Jardine's argument \cite{Jardine} and it is moreover
more or less in this form that the second author copied
Jardine's argument in \cite{nCAT}. On the other hand, the argument is found
(scattered) in \cite{JardineGoerssBook}, or else (with
an analogous list of axioms) in \cite{JardineGoerssLoc}. See also
Hovey's statement \cite{HoveyBook} 2.1.11. We advise the reader who
wants to understand this proof to refer to the references
\cite{Hirschhorn} \cite{DHK} \cite{Jardine} \cite{JardineGoerssBook}
\cite{JardineGoerssLoc} \cite{HoveyBook}.

The following lemma is therefore found
entirely in the above references.
\begin{lemma}
\label{dhklemme}
Let $M$ be a category with two subcategories $W\subset M$ and
$cof\subset M$
said to be of weak equivalences and of cofibrations (and $cof \cap W$
is said to be the
subcategory of trivial cofibrations). Suppose that:
\newline
(0) \,\, $M$ admits all limits and colimits indexed by small
categories;
\newline
(1) \,\, every retract of a morphism of $W$ (resp.\ of $cof$)
is in $W$ (resp.\ in $cof$);
\newline
(2)\,\, $W$ satisfies ``three for the price of two'': if two among $f$, $g$,
$fg$ are in $W$ then so is the third;
\newline
(3)\,\, every morphism which possesses the
lifting property with respect to the morphisms of
$cof$, is in $W$;
\newline
(4)\,\, $cof$ admits a generating subset $I\subset cof$ permitting
the small object argument;
\newline
(5)\,\, likewise, $cof \cap W$ admits a generating subset
$J\subset cof\cap W$
permitting the small object argument;
\newline
(6)\,\, $cof$ is stable under coproduct (i.e.\ if $A\rightarrow B$ is
in $cof$ and $A\rightarrow A'$ is arbitrary, then $A'\rightarrow B\cup ^AA'$
is in $cof$) and under transfinite sequential colimit (see
\cite{Hirschhorn} \cite{DHK} for the definition); and
\newline
(7)\,\, likewise, $cof \cap W$ is stable under coproduct and transfinite
sequential
colimit.

Then, denoting by $fib\subset M$ the subcategory said ``of fibrations''
equal to the class of morphisms which possess the
lifting property with respect to
trivial cofibrations, $(M,W,cof, fib)$ is a cofibrantly generated closed model
category.
\end{lemma}
{\em Proof:}
We apply the recognition criterion (\cite{DHK} 8.1) (whose notations we adopt
for the present proof). The fact that $I$ generates $cof$ in our
sense (4), implies $cof \subset I$-${\bf cof}$.  In the other direction, by
(\cite{DHK} 7.3 (ii)), the morphisms of  $I$-${\bf cof}$ are retracts of
sequential limits of coproducts by morphisms of $I$; hence by our
conditions (1) for $cof$ and (6), one has $I$-${\bf cof}\subset cof$ and hence
$
I$-${\bf cof}= cof.
$
Likewise, using our properties (5) and (1) for $cof$ $+$ (7), one has
$
J$-${\bf cof}= cof \cap W.
$
This gives us the first conditions of \cite{DHK} 8.1 (ii) and (iii),
i.e.\ that $J$-${\bf cof}=I$-${\bf cof}\cap W$.  The conditions of the beginning of
(\cite{DHK}
8.1) are our (1) for $W$ and our (2). The condition  (\cite{DHK} 8.1
(i)) results from
(4) and (5).  Finally our (3) gives $I$-${\bf inj} \subset W$,
and as lifting for $I$ (resp.\ $J$) is equivalent to lifting
for $cof$ (resp.\ $cof \cap W$), it is clear that one has
$I$-${\bf inj} \subset J$-${\bf inj}$. This gives the second part of
(\cite{DHK} 8.1 (ii)), namely $I$-${\bf inj} \subset J$-${\bf inj} \cap W$.
\eop

{\em Remark:} It results immediately from the definition of
\cite{Hirschhorn}
\cite{DHK} that if $(M, W, cof)$ is a cofibrantly generated closed model
category, then properties (0)--(7) are satisfied.
It may be useful to note that the part of the definition of a cofibrantly
generated cmc
which goes beyond
Quillen's axioms, consists in: (0) for the infinite limits and colimits; (4)
and (5); and
(6) and (7) for the sequential limits.


{\em Remark:} When we apply this lemma, we shall not verify
conditions (4) and (5), which will be consequences of the existence of an
``underlying set'' functor and of the fact that $cof$ and $W$ will be
(set-theoretically)
reasonable.
See in Jardine \cite{Jardine} the technique necessary to obtain a rigorous
proof
of (4) and (5) in our type of situation (he treats the case of simplicial
presheaves
but his technique is very general).
On the other hand, the cofibrations will in
general simply be the morphisms which induce injections on the
underlying sets; hence (1) for $cof$ and (6) will be obvious.
The existence of limits and colimits will be an immediate consequence of the
definition of the category $M$ (generally as a category of
presheaves of sets on a certain category).  In sum, we shall
verify only (1) for $W$, (2), (3) and (7). The main difficulty will be
constituted by (7).



\subnumero{Proof of \texorpdfstring{Theorem~\ref{SeCmf}}{Theorem}}

One could simply copy the proof of \cite{nCAT}
replacing ``$n$-precategory'' by ``Segal $n$-precategory''. We shall
however indicate how to obtain the proof using the lemma
\ref{dhklemme} and using directly a part of the results of
\cite{nCAT}, without copying everything. In passing, the present
proof improves that of \S 6 of \cite{nCAT} by replacing
the {\em ad hoc} argument for condition  CM5(2) by condition (3) of lemma
\ref{dhklemme} (plus, implicitly, via this lemma, the small object argument).

As in the preceding remark, conditions (0), (1) for $cof$,
and (6) are immediate. We shall not treat conditions (4) and (5)
which are of a set-theoretic nature. It therefore remains to treat (1) for
$W$, (2), (3), and (7). Let us fix $n$ and treat the category $nSePC$ of
Segal $n$-precategories.

Let $f:A\rightarrow B$ be a morphism of Segal $n$-precategories.
Suppose that $f$ is a retract of a weak equivalence $g$.  Then
$SeCat(f)$ is a retract of $SeCat(g)$ and for every $m$, $\Pi _m \circ
SeCat(f)$ is a retract of $\Pi _m \circ SeCat(g)$. The latter are
morphisms of $(n+m)$-categories, hence the closure under retracts of the
class of equivalences between $(n+m)$-categories  (\cite{nCAT} \S 6) implies
that $\Pi _m \circ SeCat(f)$ is an equivalence. This fact being true
for every $m$, it
follows that $SeCat(f)$ is an equivalence hence, by definition, $f$ is
a weak equivalence. This proves (1) for $W$.

In the same way property (2) ``three for the price of two''
is an immediate consequence (via the construction $\Pi _m\circ SeCat$)
of the same property for $(n+m)$-categories for every $m$
(\cite{nCAT} Lemma 3.8).

For property (3),
if $n=0$, property (3) is well known for
simplicial sets; we therefore suppose $n\geq 1$, and we suppose that (3)
is known for $n-1$.  Suppose that a morphism $f:A\rightarrow B$
in $nSePC$ satisfies the lifting property with respect to every
cofibration $E\hookrightarrow E'$. With the cofibration $\emptyset
\hookrightarrow \ast$ one
obtains that $f$ is surjective on objects. We now recall a
construction of \cite{limits} which will be used often throughout the
present paper. If $E$ is an $(n-1)$-precategory one defines the $n$-precategory
$\Upsilon (E)$ with two objects $\Upsilon (E)_0 = \{ 0,1\}$ and with
$$
\Upsilon (E)_{p/}(x_0,\ldots , x_p)
$$
equal to $E$
if $x_0=\ldots = x_i=0$ and $x_{i+1}=\ldots = x_p = 1$ ($0\leq i<p$); equal
to $\ast$
if $x_0=\ldots = x_p=0$ or $x_0=\ldots = x_p=1$; and equal
to $\emptyset$
otherwise. A morphism $\Upsilon (E)\rightarrow A$
 is nothing other than
a pair of objects $(x,y)\in A_0\times A_0$ equipped with a morphism
$E\rightarrow
A_{1/}(x,y)$. We return to the proof of property (3). The
lifting property for $f$ with respect to every cofibration
applies in particular
to $\Upsilon (E)\rightarrow \Upsilon (E')$ for
every cofibration $E\rightarrow E'$ of Segal $(n-1)$-precategories. This
implies
that if $x,y\in A_0$ then the morphism
$$
A_{1/}(x,y)\rightarrow B_{1/}(x,y)
$$
of Segal $(n-1)$-precategories has the lifting property for every
cofibration of Segal $(n-1)$-precategories. Property (3) in degree $n-1$
implies that $A_{1/} (x,y)\rightarrow B_{1/}(x,y)$ is a weak
equivalence. The same type of argument (using the variant of
the construction $\Upsilon (E)$ denoted $[p]E$
in \cite{limits} 2.4.5) shows that
$$
A_{p/}(x_0,\ldots , x_p)\rightarrow B_{p/}(f(x_0),\ldots , f(x_p))
$$
is a
weak equivalence. It results clearly from the description of $SeCat$ as
composite of operations of the form $Fix$ and $Gen[m]$ (see the remarks
following lemma \ref{secatunique}), that $SeCat(A)_{1/}(x,y)\rightarrow
SeCat(B)_{1/}(f(x), f(y))$ is an equivalence. Hence $SeCat(f)$ is fully
faithful, hence it is an equivalence, and $f$ is a weak equivalence.
This gives property (3).

For property (7), let
$$
B\leftarrow     A\rightarrow C
$$
be a diagram where $A\rightarrow B$ is a trivial cofibration. As remarked
above,  the Seifert-Van Kampen theorem, Theorem 9.1 of
\cite{nCAT}, gives an equivalence
$$
(\Pi _m\circ B) \cup ^{(\Pi _m \circ A)}(\Pi _m \circ C)
\stackrel{\cong}{\rightarrow} \Pi _m\circ (B\cup ^AC).
$$
The property of preservation of trivial cofibrations by coproduct
in the framework
of (non-Segal) $(n+m)$-precategories (\cite{nCAT} Lemma 3.2) therefore implies
that
$$
\Pi _m \circ C\rightarrow
\Pi _m \circ (B\cup ^AC)
$$
is a weak equivalence. This fact being true for every $m$, it follows
that $C\rightarrow
B\cup ^AC$ is a weak equivalence (it is appropriate to recall here that by
convention we have precomposed the operation $\Pi_m$ by $SeCat$). This
gives (7).
\eop


\subnumero{Homotopy classes of morphisms}

We recall (Quillen \cite{Quillen}) that if $M$ is a closed model
category and if $W\subset M$ is the subcategory whose morphisms are
the weak equivalences, then the localised category
$$
Ho(M):= W^{-1}M
$$
in the sense of Gabriel-Zisman \cite{GabrielZisman} possesses other useful
descriptions. Let $M_c$, $M_f$ and $M_{cf}$ be the categories of cofibrant,
fibrant, and cofibrant and fibrant objects, respectively. Let $W_c$, $W_f$ and
$W_{cf}$ be the intersections of $W$ with these. Then the localisations
$W_c^{-1} M_c$, $W_f^{-1}M_f$ and $W_{cf}^{-1}M_{cf}$ are all three
equivalent (via
the functors induced by the inclusions) to $Ho(M)$.
If one denotes by $\pi M_{cf}$ the category whose objects are the objects of
$M_{cf}$ and whose sets of morphisms are the sets of {\em homotopy
classes} of morphisms of $M_{cf}$ \cite{Quillen}, then $\pi M_{cf}$
is also
equivalent to $Ho(M)$. More generally, for $A,B\in M$, if one chooses
weak equivalences $A'\cong A$ and $B\cong B'$ with $A'$ cofibrant and
$B'$ fibrant, then $Hom _{Ho(M)}(A,B)$ is the set of homotopy classes
of morphisms $A'\rightarrow B'$. Following Quillen \cite{Quillen}, we shall
adopt the notation
$$
[A,B] := Hom _{Ho(M)}(A,B)
$$
for the set of morphisms from $A$ to $B$ in $Ho(M)$, with $A,B\in
Ob(M)$.
In summary, if $A$ is cofibrant and $B$ fibrant then every element of
$[A,B]$ comes from a morphism $A\rightarrow B$ of $M$, and two such morphisms
give the same element of $[A,B]$ if and only if they are homotopic
in the sense of \cite{Quillen}.

We apply this with $M= nSePC$. In this case, one has an object
$\overline{I}$, the category with two objects $0,1$ linked by an
isomorphism, regarded as a Segal $n$-category (see the heading
``Induction'' above). This object serves as an {\em interval}, i.e.\ for every
object
$A\in nSePC$, $A\times \overline{I}$ is an object ``$A\times I$'' in the sense of
Quillen \cite{Quillen} and for $B$ fibrant, two morphisms $f,g: A\rightarrow
B$ are homotopic if and only if there exists a morphism $h: A\times
\overline{I}\rightarrow B$ with $h|_{A\times \{ 0\} }=f$ and $h|_{A\times \{ 1\}
}=g$. If $A$ and $B$ are two objects of $nSePC$ then one can compute
$$
[A,B]= Hom _{Ho(nSePC)}(A,B)
$$
by choosing a weak
equivalence $B\rightarrow B'$ to a fibrant object and (thanks to the fact that
$A$ is
automatically cofibrant) by taking the set of morphisms $A\rightarrow
B'$ modulo the homotopy relation set out above with the help of
$\overline{I}$.
%%%%%%%%
See Cordier \cite{Cordier} who makes this observation in the framework of
homotopy-coherent diagrams.
%%%%%%%% above added in v3 %%%%%%%%%
